\SourcePage{120}
\SGAchapter{XV}{无上同调态射}{Morphismes acycliques}
\begin{center}M. Artin\end{center}
\textbf{引言。}
设 $g:X\to Y$ 为\Term{scheme}{概形}的\Term{morphism}{态射}。在第一节中，我们研究使 $g$ 成为\Term{acyclic}{无上同调}\Term{morphism}{态射}的条件；也就是说，对 $Y$ 上每个\Term{torsion-sheaf}{挠层} $F$，都有 $H^q(Y,F)\simeq H^q(X,g^*F)$，且把 $Y$ 换成任意\Term{etale}{平展}于 $Y$ 的 $Y'$ 后，这仍然成立。\Term{acyclic}{无上同调}性的一个自然必要条件是\Term{geometric-fiber}{几何纤维}均\Term{acyclic}{无上同调}，即对常值\Term{torsion-sheaf}{挠层}有平凡的\Term{cohomology}{上同调}。我们将看到，此条件与一个称为 $f$ 的\Term{locally-acyclic}{局部无上同调性}的\Term{local-scheme}{局部}条件（1.11）合在一起，便蕴含 $f$ \Term{acyclic}{无上同调}。

基本定理是：\Term{smooth}{光滑}\Term{morphism}{态射}对于阶数与\Term{residue-characteristic}{剩余特征}互素的\Term{torsion-sheaf}{挠层}是\Term{locally-acyclic}{局部无上同调}的（2.1）。由此将推出\Term{smooth}{光滑}\Term{morphism}{态射}的\Term{base-change}{换基}定理（第 XVI 讲 1.1）；下一讲将阐述该定理及其最初的推论。

本讲不依赖于第 XI 至 XIV 讲，尤其不依赖于“\Term{proper}{固有}\Term{morphism}{态射}的\Term{base-change}{换基}定理”。从下一讲起，后一定理将再次发挥关键\Term{action}{作用}，但将与本讲的基本定理结合使用。

\SourcePage{121}
\SGAsection{1}{整体及局部无上同调态射的一般性质}
\begin{SGAresult}{命题}{1.1}
设 $g:X\to Y$ 为\Term{scheme}{概形}的\Term{morphism}{态射}。下列断言 (i) 至 (iii) \Term{equivalence}{等价}：
\begin{enumerate}
\item[(i)] 对每个\Term{etale}{平展}\Term{morphism}{态射} $Y'\to Y$（可取其在 $Y$ 的某个\Term{affine}{仿射}\Term{open-subset}{开集}上\Term{affine}{仿射}）及 $Y'$ 上每个\Term{set}{集合}\Term{sheaf}{层} $F$，\Term{canonical}{典范}\Term{morphism}{态射} $F\to g'_*g'^*F$ 是\Term{injectivity}{单射}（相应地，\Term{bijection}{双射}）。
\item[(ii)] 对每个这样的 $Y'\to Y$ 和 $F$，\Term{canonical}{典范}\Term{morphism}{态射} $H^0(Y',F)\to H^0(X',g'^*F)$ 是\Term{injectivity}{单射}（相应地，\Term{bijection}{双射}）。
\item[(iii)] 对每个\Term{locally-quasi-finite}{局部拟有限}\Term{morphism}{态射} $Y'\to Y$ 及 $Y'$ 上每个\Term{set}{集合}\Term{sheaf}{层} $F$，\Term{canonical}{典范}\Term{morphism}{态射} $H^0(Y',F)\to H^0(X',g'^*F)$ 是\Term{injectivity}{单射}（相应地，\Term{bijection}{双射}）。
\item[(ii bis)] （若 $f$ \Term{quasi-compact}{拟紧}且\Term{quasi-separated}{拟分离}。）与 (ii) 相同，但只取形如 $I_{Y'}$ 的常值\Term{sheaf}{层} $F$，其中 $I$ 是预先给定的\Term{set}{集合}，$\operatorname{card}I\geq2$。
\end{enumerate}
\end{SGAresult}
\begin{proof}
(i) 与 (ii) 的\Term{equivalence}{等价}性由定义立即得出，且 (iii)$\Rightarrow$(ii)$\Rightarrow$(ii bis) 显然。又有 (ii bis)$\Rightarrow$(ii)，因为可假设 $Y$ \Term{affine}{仿射}，从而 $X,Y$ 均\Term{quasi-compact}{拟紧}且\Term{quasi-separated}{拟分离}，于是可应用第 XII 讲 6.5 (i)。为证明 (i)$\Rightarrow$(iii)，先记下如下引理。
\end{proof}
\begin{SGAresult}{引理}{1.2}
设 $f:Y'\to Y$ 为\Term{locally-quasi-finite}{局部拟有限}\Term{morphism}{态射}。则 $f$“对\Term{etale-topology}{平展拓扑}在 $Y'$ 上\Term{locally-finite}{局部有限}”，即存在\Term{commutative-diagram}{交换图}
\[
\begin{tikzcd}
Y'_i\arrow[r,"a_i"]\arrow[d,"b_i"']&Y'\arrow[d,"f"]\\
Y_i\arrow[r,"c_i"']&Y,
\end{tikzcd}\tag{*}
\]
其中 $a_i,c_i$ 均\Term{etale}{平展}，$b_i$ 均\Term{finite}{有限}，且 $a_i$ 构成 $Y'$ 的一个\Term{covering}{覆盖}。\footnote{参见 EGA IV 18.12.1。}
\end{SGAresult}
\begin{proof}
设 $y'\in Y'$，$y=f(y')$。只需找到图 (*)，使 $y'$ 在 $Y'_i$ 的像中。设 $\widetilde Y$ 为 $\mathscr O_{Y,y}$ 的\Term{henselization}{Hensel 化}之谱（第 VIII 讲 4.1），$Z$ 为 $\widetilde Y'=Y'\times_Y\widetilde Y$ 在位于 $y'$ 及 $\widetilde Y$ 的\Term{closed-point}{闭点} $\widetilde y$ 上方的唯一点处的普通\Term{localization}{局部化}。由第 VIII 讲 4.1 (ii)，$Z\to\widetilde Y$ \Term{finite}{有限}，且 $Z$ 为 $\widetilde Y'$ 的诱导\Term{open-subset}{开集}。由于 $\widetilde Y$ 是在 $Y$ 上\Term{etale}{平展}的\Term{prescheme}{预概形} $Y_i$ 的\Term{limit}{极限}，立即由 EGA IV 9 可知，对适当的 $Y_i$，存在 $Y'\times_Y Y_i$ 的\Term{open-subset}{开集} $Z_i$，使 $Z=Z_i\times_{Y_i}\widetilde Y$；于是在 (*) 中取 $Y'_i=Z_i$ 即可。
\end{proof}

\SourcePage{122}
\begin{proof}[命题 1.1 证明续]
现假设 1.1 (i) 成立。显然，1.1 (iii) \Term{equivalence}{等价}于把 (i) 中的\Term{etale}{平展}\Term{morphism}{态射} $Y'\to Y$ 换成\Term{locally-quasi-finite}{局部拟有限}\Term{morphism}{态射}所得的断言。对于给定的 $Y,Y'$，这一\Term{form}{形式}的断言对\Term{etale-topology}{平展拓扑}在 $Y$ 和 $Y'$ 上均为\Term{local-scheme}{局部}的，故由引理归结到 $f:Y'\to Y$ \Term{finite}{有限}的情形。设 $f':X'\to X$ 为 $f$ 沿 $g$ \Term{base-change}{换基}所得的\Term{morphism}{态射}。则\Term{base-change}{换基}\Term{morphism}{态射} $g^*f_*F\to f'_*g'^*F$ 是\Term{bijection}{双射}（第 VIII 讲 5.6）。由 (i)，$f_*F\to g_*g^*f_*F$ 是\Term{injectivity}{单射}（相应地，\Term{bijection}{双射}）。因而 (iii) 中的\Term{morphism}{态射} $\varepsilon:H^0(Y',F)\to H^0(Y',g'_*g'^*F)$ 是下列复合：
\[
\begin{aligned}
H^0(Y',F)&=H^0(Y,f_*F)\xrightarrow{\varepsilon'}H^0(Y,g_*g^*f_*F)\\
&=H^0(X,g^*f_*F)\simeq H^0(X,f'_*g'^*F)=H^0(X',g'^*F).
\end{aligned}
\]
既然 $\varepsilon'$ 是\Term{injectivity}{单射}（相应地，\Term{bijection}{双射}），$\varepsilon$ 也如此。这就完成了 1.1 的证明。
\end{proof}
\begin{SGAresult}{定义}{1.3}
具有 1.1 中\Term{equivalence}{等价}性质 (i) 至 (iii) 之一（相应地，各条\Term{lie-bracket}{括号}中的性质）的\Term{morphism}{态射} $g:X\to Y$，称为\Term{acyclic}{$(-1)$-无上同调}\Term{morphism}{态射}（相应地，\Term{acyclic}{$0$-无上同调}\Term{morphism}{态射}）。若对每个\Term{base-change}{换基}\Term{morphism}{态射} $Y'\to Y$，\Term{morphism}{态射} $g'=g\times_Y Y':X\times_Y Y'\to Y'$ 均为\Term{acyclic}{$(-1)$-无上同调}（相应地，\Term{acyclic}{$0$-无上同调}），则称原\Term{morphism}{态射}\Term{universally-acyclic}{普适 $(-1)$-无上同调}（相应地，\Term{universally-acyclic}{普适 $0$-无上同调}）。最后，称\Term{scheme}{概形} $X$ \Term{acyclic}{$(-1)$-无上同调}（相应地，\Term{acyclic}{$0$-无上同调}），若它\Term{nonempty}{非空}（相应地，\Term{connected}{连通}且\Term{nonempty}{非空}）。
\end{SGAresult}
\begin{SGAresult}{推论}{1.4}
设 $g:X\to Y$ 为\Term{acyclic}{$0$-无上同调}\Term{morphism}{态射}。则对每个\Term{locally-quasi-finite}{局部拟有限}\Term{morphism}{态射} $Y'\to Y$ 及 $Y'$ 上每个群\Term{sheaf}{层} $F$，\Term{canonical}{典范}\Term{morphism}{态射}
\[H^1(Y',F)\longrightarrow H^1(X',g'^*F)\]
是\Term{injectivity}{单射}。
\end{SGAresult}
这立即由第 XII 讲 6.5 (i) 得出。
\begin{SGAresult}{推论}{1.5}
设 $g:X\to Y$ 为\Term{morphism}{态射}。若 $g$ \Term{surjectivity}{满射}（即其\Term{fiber}{纤维}均\Term{nonempty}{非空}，亦即其\Term{fiber}{纤维}均\Term{acyclic}{$(-1)$-无上同调}），则 $g$ \Term{acyclic}{$(-1)$-无上同调}。当 $g$ \Term{quasi-compact}{拟紧}且\Term{quasi-separated}{拟分离}时，逆命题成立。
\end{SGAresult}
第一个断言采用 1.1 (ii) 的\Term{form}{形式}并考察所涉及的\Term{sheaf}{层}的茎即显然。第二个断言由定义的 1.1 (iii) \Term{form}{形式}和第 VII 讲 5 的取\Term{limit}{极限}法则得出；这里注意，对任意 $y\in Y$，$y=\operatorname{Spec}k(y)$ 是在 $Y$ 上\Term{locally-quasi-finite}{局部拟有限}的\Term{affine}{仿射}\Term{scheme}{概形}的\Term{inverse-limit}{投射极限}，即 $y$ 在赋予诱导\Term{reduced-structure}{约化结构}的 $\overline{\{y\}}$ 中的\Term{affine}{仿射}\Term{neighborhood}{邻域}的\Term{inverse-limit}{投射极限}。注意，1.5 蕴含：当 $g$ \Term{quasi-compact}{拟紧}且\Term{quasi-separated}{拟分离}时，$g$ \Term{acyclic}{$(-1)$-无上同调}已蕴含它\Term{universally-acyclic}{普适 $(-1)$-无上同调}。
\begin{SGAresult}{命题}{1.6}
设 $g:X\to Y$ 为\Term{scheme}{概形}的\Term{morphism}{态射}，$L\subset\mathbb P$ 为一个素数集，$n\in\mathbb N$ 为自然数（相应地，$n=1$）。下列断言 (i) 至 (iii)（相应地，各条\Term{lie-bracket}{括号}中的断言）\Term{equivalence}{等价}。
\SourcePage{123}
若进一步 $g$ \Term{quasi-compact}{拟紧}且\Term{quasi-separated}{拟分离}，则这些断言还\Term{equivalence}{等价}于 (iv)（相应地，其\Term{lie-bracket}{括号}中的\Term{form}{形式}）。
\begin{enumerate}
\item[(i)] 对每个\Term{etale}{平展}\Term{morphism}{态射} $Y'\to Y$ 及 $Y'$ 上每个阿贝尔 $L$-\Term{torsion-sheaf}{挠层}（相应地，$\operatorname{ind}$-$L$-群\Term{sheaf}{层}，第 IX 讲 1.5）$F$，有
\[
F\xrightarrow{\sim}g'_*g'^*F,\qquad (R^qg'_*)g'^*F=0\quad(1\leq q\leq n\text{，相应地 }q=1).
\]
\item[(ii)] 对每个\Term{etale}{平展}\Term{morphism}{态射} $Y'\to Y$ 及 $Y'$ 上每个阿贝尔 $L$-\Term{torsion-sheaf}{挠层}（相应地，$\operatorname{ind}$-$L$-群\Term{sheaf}{层}）$F$，\Term{canonical}{典范}\Term{morphism}{态射}
\[H^q(Y',F)\longrightarrow H^q(X',g'^*F)\]
在 $q\leq n$ 时\Term{bijection}{双射}，在 $q=n+1$ 时\Term{injectivity}{单射}（相应地，在 $q\leq1$ 时\Term{bijection}{双射}）。
\item[(iii)] 对每个\Term{locally-quasi-finite}{局部拟有限}\Term{morphism}{态射} $Y'\to Y$ 及 $Y'$ 上每个阿贝尔 $L$-\Term{torsion-sheaf}{挠层}（相应地，$\operatorname{ind}$-$L$-群\Term{sheaf}{层}）$F$，\Term{morphism}{态射}
\[H^q(Y',F)\longrightarrow H^q(X',g'^*F)\]
在 $q\leq n$ 时\Term{bijection}{双射}，在 $q=n+1$ 时\Term{injectivity}{单射}（相应地，在 $q\leq1$ 时\Term{bijection}{双射}）。
\item[(iv)] \Term{morphism}{态射} $g$ \Term{surjectivity}{满射}，且对每个\Term{locally-quasi-finite}{局部拟有限}\Term{morphism}{态射} $Y'\to Y$、每个 $\nu>0$ 和每个 $\ell\in L$，\Term{canonical}{典范}\Term{morphism}{态射}
\[H^q(Y',\mathbb Z/\ell^\nu)\longrightarrow H^q(X',\mathbb Z/\ell^\nu)\]
对 $0\leq q\leq n$ \Term{surjectivity}{满射}（相应地，对每个\Term{locally-quasi-finite}{局部拟有限}\Term{morphism}{态射} $Y'\to Y$ 及每个普通\Term{finite}{有限} $L$-群 $G$，\Term{canonical}{典范}\Term{morphism}{态射} $H^i(Y',G)\to H^i(X',G)$ 对 $i=0,1$ \Term{surjectivity}{满射}）。
\end{enumerate}
\end{SGAresult}
\begin{proof}
(ii)$\Leftarrow$(iii)$\Rightarrow$(iv) 显然。又 $(R^qg'_*)g'^*F$ 是\Term{presheaf}{预层} $R^q(Y'')=H^q(X'',g''^*F)$（$Y''\to Y'$ \Term{etale}{平展}）的\Term{adjunction}{伴随}\Term{sheaf}{层}。若 (ii) 成立，则 $H^q(Y'',F)\xrightarrow{\sim}H^q(X'',g''^*F)$，故此\Term{adjunction}{伴随}\Term{sheaf}{层}为零；即 (ii)$\Rightarrow$(i)。(i)$\Rightarrow$(ii) 在阿贝尔情形由 1.1 (i)（对应 $n=0$）与 Leray \Term{spectral-sequence}{谱序列}
\[H^p(Y',(R^qg'_*)g'^*F)\Longrightarrow H^{p+q}(X',g'^*F)\]
得出；在群\Term{sheaf}{层}情形，则使用\Term{exact-sequence}{正合列}
\[0\longrightarrow H^1(Y',g'_*g'^*F)\longrightarrow H^1(X',g'^*F)\longrightarrow H^0(Y',(R^1g'_*)g'^*F).\]
(i)$\Rightarrow$(iii)：由于 (i) 与 (ii) \Term{equivalence}{等价}，显然 (iii) \Term{equivalence}{等价}于将 (i) 中“\Term{etale}{平展}”一词替换为“\Term{locally-quasi-finite}{局部拟有限}”所得的断言。对给定的 $Y,Y'$，这一\Term{form}{形式}的断言对\Term{etale-topology}{平展拓扑}在 $Y'$ 和 $Y$ 上均为\Term{local-scheme}{局部}的，故由 1.2 归结到 $f:Y'\to Y$ \Term{finite}{有限}的情形。
\SourcePage{124}
设 $f':X'\to X$ 为相应\Term{morphism}{态射}。三次应用第 VIII 讲 5.5（相应地，第 VIII 讲 5.8），有
\[
\begin{aligned}
f_*(R^qg'_*)g'^*F&\simeq R^q(fg')_*g'^*F\simeq R^q(gf')_*g'^*F\\
&\simeq(R^qg_*)f'_*g'^*F\simeq(R^qg_*)g^*(f_*F).
\end{aligned}
\]
若 (i) 成立，则最后一项对所考虑的 $q$ 为零，故 (i)$\Rightarrow$(iii)。最后还须证明 (iv)$\Rightarrow$(iii)，而这立即由第 XII 讲 6.5 得出。
\end{proof}
\begin{SGAresult}{定义}{1.7}
设 $n\in\mathbb N$，$L\subset\mathbb P$。满足 1.6 中\Term{equivalence}{等价}条件 (i) 至 (iii) 之一（相应地，各条\Term{lie-bracket}{括号}中条件之一）的\Term{morphism}{态射} $g:X\to Y$，称为对 $L$ \Term{acyclic}{$n$-无上同调}的\Term{morphism}{态射}（相应地，对 $L$ \Term{one-aspherical}{$1$-非球面}的\Term{morphism}{态射}）。若 $g$ 对每个 $n$ 均对 $L$ \Term{acyclic}{$n$-无上同调}，则称其对 $L$ \Term{acyclic}{无上同调}。若对每个 $Y'\to Y$，\Term{morphism}{态射} $g'=g\times_Y Y'$ 均对 $L$ \Term{acyclic}{$n$-无上同调}（相应地，\Term{one-aspherical}{$1$-非球面}），则称 $g$ 对 $L$ \Term{universally-acyclic}{普适 $n$-无上同调}（相应地，\Term{universally-one-aspherical}{普适 $1$-非球面}）。最后，称\Term{prescheme}{预概形} $X$ 对 $L$ \Term{acyclic}{$n$-无上同调}（相应地，\Term{one-aspherical}{$1$-非球面}），若它 \Term{acyclic}{$0$-无上同调}（1.3），且对每个 $\ell\in L$ 及 $1\leq q\leq n$，$H^q(X,\mathbb Z/\ell)=0$（相应地，对每个普通\Term{finite}{有限} $L$-群 $G$，$H^1(X,G)=0$）。
\end{SGAresult}
\textbf{注记 1.8。}
a) 若 $g$ \Term{quasi-compact}{拟紧}且\Term{quasi-separated}{拟分离}，则为使 $g$ \Term{universally-acyclic}{普适 $n$-无上同调}，只需每当 $Y'\to Y$ \Term{locally-finitely-presented}{局部有限呈示}时，$g'$ 均\Term{acyclic}{$n$-无上同调}。我们不需要这一事实，将其证明留给读者作为取\Term{limit}{极限}的练习。

b) 我们不知道对 $L$ \Term{acyclic}{$n$-无上同调}的\Term{morphism}{态射}是否已对 $L$ \Term{universally-acyclic}{普适 $n$-无上同调}。对于下文 1.11 的\Term{local-scheme}{局部}变体，也有同样问题。

c) 当 $n=0$ 时，便得到 1.3 的概念。当然，可以把 1.6、1.7 写成包含 $n=-1$ 的\Term{form}{形式}。下面的命题在此情形也成立。
\begin{SGAresult}{命题}{1.9}
(i) 设 $X\xrightarrow{g}Y\xrightarrow{h}Z$ 为\Term{scheme}{概形}的\Term{morphism}{态射}，$L\subset\mathbb P$，$n\in\mathbb N$。假设 $g$ \Term{acyclic}{$0$-无上同调}（相应地，对 $L$ \Term{one-aspherical}{$1$-非球面}，或对 $L$ \Term{acyclic}{$n$-无上同调}）。令 $f=hg$。则对 $Z$ 上任意\Term{set}{集合}\Term{sheaf}{层}（相应地，$\operatorname{ind}$-$L$-群\Term{sheaf}{层}，或阿贝尔 $L$-\Term{torsion-sheaf}{挠层}）$F$，有\Term{isomorphism}{同构} $f_*f^*F\xrightarrow{\sim}h_*h^*F$（相应地，$p\leq1$，或 $p\leq n$ 时，有\Term{isomorphism}{同构} $(R^pf_*)f^*F\xrightarrow{\sim}(R^ph_*)h^*F$）。特别地，$h$ \Term{acyclic}{$0$-无上同调}（相应地，对 $L$ \Term{one-aspherical}{$1$-非球面}，或对 $L$ \Term{acyclic}{$n$-无上同调}）当且仅当 $f$ 如此。

(ii) 设 $\{g_\alpha:X_\alpha\to Y_\alpha\}$ 为\Term{scheme}{概形}\Term{morphism}{态射}的\Term{inverse-system}{投射系}，其中 $X_\alpha\to X_\beta$ 及 $Y_\alpha\to Y_\beta$ 均\Term{affine}{仿射}，且各 $g_\alpha$ \Term{quasi-compact}{拟紧}且\Term{quasi-separated}{拟分离}。令 $g:X\to Y$ 为各 $g_\alpha$ 的\Term{inverse-limit}{投射极限}（第 VII 讲 5.1）。若所有 $g_\alpha$ 均\Term{acyclic}{$0$-无上同调}（相应地，对 $L$ \Term{one-aspherical}{$1$-非球面}，或对 $L$ \Term{acyclic}{$n$-无上同调}），则 $g$ 也如此。
\end{SGAresult}
\SourcePage{125}
\begin{proof}
对于 (i)，首先有
\[h_*h^*F\xrightarrow{\sim}h_*(g_*g^*)h^*F\simeq f_*f^*F.\]
(i) 的第二个断言是 Leray \Term{spectral-sequence}{谱序列}
\[(R^ph_*)(R^qg_*)f^*F\Longrightarrow(R^{p+q}f_*)f^*F\]
的直接推论。最后，非阿贝尔断言由第 XII 讲 3.2 的\Term{exact-sequence}{正合列}得出。(ii) 容易由第 VII 讲 5 及第 VI 讲的取\Term{limit}{极限}理论得出。
\end{proof}
\begin{SGAresult}{推论}{1.9.1}
设 $g:X\to Y$ 为\Term{quasi-compact}{拟紧}且\Term{quasi-separated}{拟分离}的\Term{morphism}{态射}，$n\in\mathbb N$，$L\subset\mathbb P$。若 $g$ 对 $L$ \Term{acyclic}{$n$-无上同调}（相应地，对 $L$ \Term{one-aspherical}{$1$-非球面}），则对 $Y$ 上任意在其所在点 $y\in Y$ 上\Term{algebra}{代数}的\Term{geometric-point}{几何点} $\bar y$，\Term{geometric-fiber}{几何纤维} $X_{\bar y}$ 对 $L$ \Term{acyclic}{$n$-无上同调}（相应地，对 $L$ \Term{one-aspherical}{$1$-非球面}）。
\end{SGAresult}
事实上，$\bar y$ 是在 $Y$ 上\Term{locally-quasi-finite}{局部拟有限}的\Term{affine}{仿射}\Term{scheme}{概形} $Y_i$ 的\Term{inverse-limit}{投射极限}；由于 $X_i=X\times_Y Y_i\to Y_i$ 对 $L$ \Term{acyclic}{$n$-无上同调}（相应地，……），由 1.9 (ii)，$X_{\bar y}\to\bar y$ 也如此，因而得到所述结论。
\begin{SGAresult}{命题}{1.10}
设 $g:X\to Y$ 为\Term{scheme}{概形}的\Term{morphism}{态射}，$L\subset\mathbb P$，$n\in\mathbb N$。下列条件 (i) 至 (iv) \Term{equivalence}{等价}。
\begin{enumerate}
\item[(i)] 设 $\bar y$ 为 $Y$ 的\Term{geometric-point}{几何点}，$\bar x$ 为其上方 $X$ 的\Term{geometric-point}{几何点}。设 $\widetilde X,\widetilde Y$ 为 $X,Y$ 在 $\bar x,\bar y$ 处的\Term{strict-localization}{严格局部化}，$\widetilde g:\widetilde X\to\widetilde Y$ 为 $g$ 诱导的\Term{morphism}{态射}。则 $\widetilde g$ \Term{acyclic}{$0$-无上同调}（相应地，对 $L$ \Term{one-aspherical}{$1$-非球面}，或对 $L$ \Term{acyclic}{$n$-无上同调}）。
\item[(ii)] 对每个方块均为\Term{cartesian}{笛卡尔}的图
\[
\begin{tikzcd}
X\arrow[d,"g"']&X'\arrow[l,"j"']\arrow[d,"g'"]&X''\arrow[l,"j'"']\arrow[d,"g''"]\\
Y&Y'\arrow[l,"i"]&Y''\arrow[l,"i'"]
\end{tikzcd}\tag{*}
\]
其中 $i$ \Term{etale}{平展}，$i'$ 为\Term{finitely-presented}{有限呈示}的\Term{etale}{平展}\Term{morphism}{态射}，以及 $Y''$ 上每个\Term{set}{集合}\Term{sheaf}{层}（相应地，$\operatorname{ind}$-$L$-群\Term{sheaf}{层}，或阿贝尔 $L$-\Term{torsion-sheaf}{挠层}）$F$，\Term{base-change}{换基}\Term{morphism}{态射}（第 XII 讲 4.1、4.2）
\[g'^*i'_*F\longrightarrow j'_*g''^*F\]
为\Term{bijection}{双射}（相应地，\Term{morphism}{态射}
\[g'^*(R^qi'_*)F\longrightarrow(R^qj'_*)g''^*F\]
对 $q=0,1$ 为\Term{bijection}{双射}；或对 $0\leq q\leq n$ 为\Term{bijection}{双射}，对 $q=n+1$ 为\Term{injectivity}{单射}）。
\item[(iii)] 与 (ii) 相同，但仅假设 $i'$ \Term{quasi-finite}{拟有限}且\Term{quasi-separated}{拟分离}。
\SourcePage{126}
\item[(iv)] 与 (ii) 相同，但仅假设 $i$ \Term{locally-quasi-finite}{局部拟有限}，而要求 $i'$ 为\Term{quasi-compact}{拟紧}的\Term{open-immersion}{开浸入}。
\end{enumerate}
\end{SGAresult}
\begin{proof}
(iii)$\Rightarrow$(ii) 显然。验证 (i)$\Rightarrow$(iii)。只需逐茎考察 (iii) 中的\Term{morphism}{态射}。设 $\bar y$ 为 $Y'$ 的\Term{geometric-point}{几何点}，$\bar x$ 为其上方 $X'$ 的\Term{geometric-point}{几何点}。设
\[
\begin{tikzcd}
\widetilde X'\arrow[d,"\widetilde g'"']&\widetilde X''\arrow[l,"\widetilde j'"']\arrow[d,"\widetilde g''"]\\
\widetilde Y'&\widetilde Y''\arrow[l,"\widetilde i'"]
\end{tikzcd}
\]
为由 (*) 得到的\Term{cartesian}{笛卡尔}图，其中 $\widetilde X',\widetilde Y'$ 是 $X',Y'$ 在 $\bar x,\bar y$ 处的\Term{strict-localization}{严格局部化}，且 $\widetilde Y''=Y''\times_{Y'}\widetilde Y'$，$\widetilde X''=X''\times_{X'}\widetilde X'$。由第 VIII 讲 5.2、5.3，
\[H^q(\widetilde Y'',\widetilde F)\simeq(R^qi'_*F)_{\bar y}\simeq(g'^*(R^qi'_*)F)_{\bar x},\]
其中 $\widetilde F$ 为 $F$ 在 $\widetilde Y''$ 上的\Term{inverse-image}{逆像}，且
\[H^q(\widetilde X'',\widetilde g''^*\widetilde F)\simeq((R^qj'_*)g''^*F)_{\bar x}.\]
由于 $\widetilde i'$ \Term{quasi-finite}{拟有限}，对所考察的 $q$ 可将 1.1 (iii)（相应地，1.5 (iii)）应用于\Term{morphism}{态射}
\[H^q(\widetilde Y',\widetilde F)\longrightarrow H^q(\widetilde X',\widetilde g''^*\widetilde F),\]
故得 (i)$\Rightarrow$(iii)。

(ii)$\Rightarrow$(i)。假设 (ii) 成立，验证\Term{morphism}{态射} $\widetilde g:\widetilde X\to\widetilde Y$ 的\Term{acyclic}{无上同调}性判据 1.6 (ii)。设 $\widetilde Y'\to\widetilde Y$ \Term{etale}{平展}，$\widetilde Y'$ \Term{affine}{仿射}，$\widetilde F$ 为 $\widetilde Y'$ 上的\Term{set}{集合}\Term{sheaf}{层}（相应地，阿贝尔 $L$-\Term{torsion-sheaf}{挠层}，或 $\operatorname{ind}$-$L$-群\Term{sheaf}{层}）。须证明
\[H^q(\widetilde Y',\widetilde F)\xrightarrow{\sim}H^q(\widetilde X',\widetilde g'^*\widetilde F),\]
其中 $\widetilde g':\widetilde X'\to\widetilde Y'$ 由 $\widetilde g$ \Term{base-change}{换基}而得。写成 $\widetilde Y=\varprojlim Y_\alpha$，其中 $Y_\alpha$ 在 $Y$ 上\Term{etale}{平展}且\Term{affine}{仿射}（可假设 $Y$ \Term{affine}{仿射}）。可以把\Term{etale}{平展}\Term{morphism}{态射} $\widetilde Y'\to\widetilde Y$“\Term{descent}{下降}”到某个 $Y_\alpha$，记为 $Y'_0\to Y_0$，且有 $\widetilde Y'=\varprojlim Y'_\alpha$，$Y'_\alpha=Y_\alpha\times_{Y_0}Y'_0$。$\widetilde Y'$ 上每个\Term{sheaf}{层} $\widetilde F$ 均为 $f_\alpha'^*f'_{\alpha*}\widetilde F$ 的\Term{inductive-limit}{归纳极限}，其中 $f'_\alpha:\widetilde Y'\to Y'_\alpha$ 为\Term{canonical}{典范}\Term{morphism}{态射}。由第 XII 讲 1.2 (v)、1.6 (iii)，$f'_{\alpha*}\widetilde F$ 为\Term{set}{集合}\Term{sheaf}{层}（相应地，阿贝尔 $L$-\Term{torsion-sheaf}{挠层}，或 $\operatorname{ind}$-$L$-群\Term{sheaf}{层}）。结合第 VII 讲 3.3，只需处理 $\widetilde F$ \Term{descent}{下降}到某个 $Y'_\alpha$，例如 $Y'_0$ 上的 $F_0$ 的情形。
\SourcePage{127}
将 (ii) 应用于图
\[
\begin{tikzcd}
X\arrow[d,"g"']&X_0\arrow[l]\arrow[d,"g_0"]&X'_0\arrow[l,"j'_0"']\arrow[d,"g'_0"]\\
Y&Y_0\arrow[l]&Y'_0\arrow[l,"i'_0"]
\end{tikzcd}
\]
便得，对适用的 $q$，$g_0^*(R^qi'_{0*})F_0\xrightarrow{\sim}(R^qj'_{0*})g_0'^*F_0$。\Term{geometric-point}{几何点} $\bar y,\bar x$ \Term{canonical}{典范}地\Term{lifting}{提升}为 $Y_0,X_0$ 的\Term{geometric-point}{几何点}，仍以相同\Term{signature}{符号}表示。由第 VIII 讲 5.2、5.3，
\[H^q(\widetilde Y',\widetilde F)=((R^qi'_{0*})F_0)_{\bar y}=(g_0^*(R^qi'_{0*})F_0)_{\bar x},\]
以及
\[H^q(\widetilde X',\widetilde g'^*\widetilde F)=((R^qj'_{0*})g_0'^*F_0)_{\bar x},\]
对所考察的 $q$ 成立，故得结论。

还须证明 (i) 至 (iii) 与 (iv) \Term{equivalence}{等价}。采用 (i) 的\Term{form}{形式}，可见这些条件在\Term{locally-quasi-finite}{局部拟有限}的底\Term{group-extension}{扩张} $Y'\to Y$ 下\Term{stable}{稳定}；由此在 (ii) 的\Term{form}{形式}下立即得到它们蕴含 (iv)。反之，(iv) 蕴含 (iii)：为验证 (iii)，例如利用 Leray \Term{spectral-sequence}{谱序列}或第 XII 讲 6.8 的\Term{descent}{下降}结果，立即可知可假设 $Y',Y''$ \Term{affine}{仿射}。于是由“主定理”（EGA IV 8.12.8），$Y''\to Y'$ \Term{resolution}{分解}为 $Y''\xrightarrow{i'_1}Y'_1\xrightarrow{i'_2}Y'$，其中 $i'_1$ 为\Term{open-immersion}{开浸入}，$i'_2$ 为\Term{finite-morphism}{有限态射}。对 $i'_2$ 使用第 VIII 讲 5.5、5.8，便归结到为 $i'_1$ 验证 (ii)，而这属于 (iv)。这就完成了 1.10 的证明。
\end{proof}
\begin{SGAresult}{定义}{1.11}
满足 1.10 中\Term{equivalence}{等价}条件 (i) 至 (iv) 之一的\Term{prescheme}{预概形}\Term{morphism}{态射} $g:X\to Y$，称为\Term{locally-acyclic}{局部 $0$-无上同调}\Term{morphism}{态射}（相应地，对 $L$ \Term{locally-one-aspherical}{局部 $1$-非球面}的\Term{morphism}{态射}，或对 $L$ \Term{locally-acyclic}{局部 $n$-无上同调}的\Term{morphism}{态射}）。以显然的方式（比较 1.7）定义对 $L$ \Term{locally-acyclic}{局部无上同调}\Term{morphism}{态射}及\Term{universally-locally-acyclic}{普适局部 $0$-无上同调}\Term{morphism}{态射}（相应地，……）的概念。
\end{SGAresult}
自然可以变通 1.10 (i) 至 (iii) 的陈述。例如，在 (iii)、(iv) 中，可以只取形如 $G_I$ 的常值\Term{sheaf}{层} $F$，其中 $G$ 为\Term{finite}{有限}\Term{set}{集合}（相应地，\Term{finite}{有限} $L$-群，或形如 $\mathbb Z/\ell^\nu\mathbb Z$，$\ell\in L$，$\nu\in\mathbb N$ 的群）；这容易由第 VIII 讲 5.2、5.3 及第 XII 讲 6.5 得出。还指出如下变体。
\begin{SGAresult}{推论}{1.12}
假设 $g:X\to Y$ \Term{locally-finite-type}{局部有限型}。为验证\Term{locally-acyclic}{局部 $0$-无上同调}性（相应地，……），只需在\Term{geometric-point}{几何点} $\bar x$ 在\Term{geometric-fiber}{几何纤维} $X_{\bar y}$ 中为\Term{closed-point}{闭点}的情形验证 1.10 (i)。
\end{SGAresult}
\SourcePage{128}
事实上，在 1.10 的 (i)$\Rightarrow$(iii) 的证明中，我们逐茎验证了 (iii)；由于 $g$ \Term{locally-finite-type}{局部有限型}，只需对在其\Term{geometric-fiber}{几何纤维}中为\Term{closed-point}{闭点}的\Term{geometric-point}{几何点} $\bar x$ 处的茎验证即可（第 VIII 讲 3.13 b)）。
\begin{SGAresult}{命题}{1.13}
(i) 设 $X\xrightarrow{g}Y\xrightarrow{h}Z$ 为\Term{scheme}{概形}的\Term{morphism}{态射}，$L\subset\mathbb P$，$n\in\mathbb N$。假设 $g$ \Term{surjectivity}{满射}，且\Term{locally-acyclic}{局部 $0$-无上同调}（相应地，对 $L$ \Term{locally-one-aspherical}{局部 $1$-非球面}，或对 $L$ \Term{locally-acyclic}{局部 $n$-无上同调}）。则 $h$ \Term{locally-acyclic}{局部 $0$-无上同调}（相应地，……），当且仅当 $f=hg$ 如此。

(ii) 设 $\{g_\alpha:X_\alpha\to Y_\alpha\}$ 为\Term{prescheme}{预概形}\Term{morphism}{态射}的\Term{inverse-system}{投射系}，其中 $X_\alpha\to X_\beta$ 和 $Y_\alpha\to Y_\beta$ 均\Term{affine}{仿射}。设 $g$ 为此系的\Term{inverse-limit}{投射极限}（第 VII 讲 5.1）。若所有 $g_\alpha$ 均\Term{locally-acyclic}{局部 $0$-无上同调}（相应地，……），则 $g$ 也如此。
\end{SGAresult}
这是与 1.9 类似的一组\Term{form}{形式}性质。

\textbf{注记 1.14。}
为简化工作，执笔人在 1.10 至 1.13 中没有同时处理 $n=-1$ 的情形；读者可自行确信，这些陈述、定义和证明基本上原样适用于该情形。例如，沿用 1.10 的记号，条件 (i) 表述为：诱导的\Term{local-morphism}{局部态射} $\widetilde X\to\widetilde Y$ 均\Term{acyclic}{$(-1)$-无上同调}，即\Term{surjectivity}{满射}（1.5）；条件 (ii) 至 (iv) 则表述为：\Term{base-change}{换基}\Term{morphism}{态射} $g'^*(i'_*(F))\to j'_*(g''^*(F))$ 为\Term{monomorphism}{单态射}。满足这些条件时，便称 $g$ \Term{locally-acyclic}{局部 $(-1)$-无上同调}。因此，此条件在\Term{locally-quasi-finite}{局部拟有限}底\Term{group-extension}{扩张}下\Term{stable}{稳定}。例如，\Term{flat}{平坦}\Term{morphism}{态射}\Term{locally-acyclic}{局部 $(-1)$-无上同调}，因为此时 $\widetilde X\to\widetilde Y$ 的\Term{surjectivity}{满射}性是这些\Term{morphism}{态射}\Term{flat}{平坦}的一个熟知推论。另一方面，当 $f$ \Term{locally-finitely-presented}{局部有限呈示}时，$f$ \Term{locally-acyclic}{局部 $(-1)$-无上同调}当且仅当它\Term{universally-open}{普适开}，至少在 $Y$ \Term{locally-noetherian}{局部 Noether}时如此；这容易由 EGA IV 14.4.9、1.10.3 得出。故对此种\Term{morphism}{态射}，\Term{locally-acyclic}{局部 $(-1)$-无上同调}性蕴含\Term{universally-locally-acyclic}{普适局部 $(-1)$-无上同调}性。（此外，上述最后几个陈述很可能\footnote{M. Raynaud 实际上已验证这一点。参见 EGA IV 14 的再版！}不需要 $Y$ \Term{locally-noetherian}{局部 Noether}的假设；若 $Y$ 的\Term{irreducible-component}{不可约分支}族\Term{locally-finite}{局部有限}，则确实如此，见 EGA IV 14.4.10。）
\begin{SGAresult}{定理}{1.15}
设 $g:X\to Y$ 为\Term{quasi-compact}{拟紧}且\Term{quasi-separated}{拟分离}的\Term{morphism}{态射}，$L\subset\mathbb P$ 为素数集，$n$ 为自然数（相应地，$n=1$）。假设 $g$ 对 $L$ \Term{acyclic}{$(n-1)$-无上同调}且对 $L$ \Term{locally-acyclic}{局部 $(n-1)$-无上同调}（1.3、1.7、1.11、1.14）。设 $F$ 为 $Y$ 上的\Term{sheaf}{层}，当 $n\geq1$ 时假设它具有阿贝尔 $L$-\Term{torsion-sheaf}{挠层}结构（相应地，群\Term{sheaf}{层}结构），$\xi\in H^n(X,g^*(F))$。则 $\xi$ 来自 $H^n(Y,F)$ 的某个元素（由 $g$ 的\Term{acyclic}{$(n-1)$-无上同调}性，此元素唯一），当且仅当
\SourcePage{129}
对 $Y$ 上每个在其所在点 $y\in Y$ 上\Term{algebra}{代数}的\Term{geometric-point}{几何点} $\bar y$，$\xi$ 在 $H^n(X_{\bar y},F|_{X_{\bar y}})$ 中的像 $\xi_{\bar y}$ 属于 $H^n(\bar y,F_{\bar y})$ 的像（当 $n\geq1$ 时，这意味着 $\xi_{\bar y}=0$）。
\end{SGAresult}
立即指出，1.15 蕴含：在上述关于 $g$ 的预备条件下，$g$ 对 $L$ \Term{acyclic}{$n$-无上同调}（相应地，对 $L$ \Term{one-aspherical}{$1$-非球面}），当且仅当对每个上述 $\bar y$，$X_{\bar y}$ 对 $L$ \Term{acyclic}{$n$-无上同调}（相应地，对 $L$ \Term{one-aspherical}{$1$-非球面}）。“若”的方向直接由 1.15 及定义得出，“仅若”的方向由 1.9.1 总是立即成立。因此，直接对 $n$ 归纳便得：
\begin{SGAresult}{推论}{1.16}
设 $g$ 为\Term{quasi-compact}{拟紧}且\Term{quasi-separated}{拟分离}的\Term{morphism}{态射}，它对 $L$ \Term{locally-acyclic}{局部 $(n-1)$-无上同调}，其中 $n$ 为给定自然数（相应地，$n=1$）。则 $g$ 对 $L$ \Term{acyclic}{$n$-无上同调}（相应地，对 $L$ \Term{one-aspherical}{$1$-非球面}），当且仅当对 $Y$ 上每个在其所在点 $y\in Y$ 上\Term{algebra}{代数}的\Term{geometric-point}{几何点} $\bar y$，\Term{fiber}{纤维} $X_{\bar y}$ 对 $L$ \Term{acyclic}{$n$-无上同调}（相应地，对 $L$ \Term{one-aspherical}{$1$-非球面}）。
\end{SGAresult}
\begin{SGAresult}{推论}{1.17}
设 $g:X\to Y$ 为\Term{prescheme}{预概形}的\Term{morphism}{态射}，$L\subset\mathbb P$ 为素数集，$n$ 为自然数（相应地，$n=1$）。则 $g$ 对 $L$ \Term{locally-acyclic}{局部 $n$-无上同调}（相应地，对 $L$ \Term{locally-one-aspherical}{局部 $1$-非球面}）的充要条件如下：对 $X$ 的每个\Term{geometric-point}{几何点} $\bar x$，若 $\widetilde g:\widetilde X\to\widetilde Y$ 是 $g$ 在 $X,Y$ 的\Term{strict-localization}{严格局部化}之间诱导的\Term{morphism}{态射}（分别在 $\bar x$ 和 $Y$ 上相应\Term{geometric-point}{几何点} $\bar y$ 处），则 $g$ 对于 $Y$ 上每个在其所在点 $z\in Y$ 上\Term{algebra}{代数}的\Term{geometric-point}{几何点} $\bar z$ 的\Term{geometric-fiber}{几何纤维} $X_{\bar z}$，均对 $L$ \Term{acyclic}{$n$-无上同调}（相应地，对 $L$ \Term{one-aspherical}{$1$-非球面}）。
\end{SGAresult}
必要性由判据 1.10 (i) 及 1.9.1 得出。为证充分性，对 $n$ 归纳；当 $n\geq1$ 时，可假设 $g$ 已对 $L$ \Term{locally-acyclic}{局部 $(n-1)$-无上同调}。于是各\Term{morphism}{态射} $\widetilde g$ 也如此，而 1.16 蕴含 $\widetilde g$ \Term{acyclic}{$n$-无上同调}（相应地，对 $L$ \Term{one-aspherical}{$1$-非球面}）；由判据 1.10 (i)，这蕴含 $g$ 对 $L$ \Term{locally-acyclic}{局部 $n$-无上同调}（相应地，\Term{locally-one-aspherical}{局部 $1$-非球面}）。
\begin{SGAresult}{推论}{1.18}
沿用 1.17 的记号，若 $g$ \Term{locally-finite-type}{局部有限型}，则在该判据中，只需取其所在点 $x$ 在相应\Term{fiber}{纤维}中为\Term{closed-point}{闭点}的\Term{geometric-point}{几何点} $\bar x$。
\end{SGAresult}
这由上述证明及 1.12 得出。
\begin{proof}[1.15 的证明\footnote{只关心 Noether 情形的读者可参阅 1.18。}]
先证明可以假设 $Y$ \Term{strictly-local}{严格局部}。在阿贝尔情形 $n\geq1$，$g$ 和 $g^*(F)$ 的 Leray \Term{spectral-sequence}{谱序列}，结合表达 $g$ 的\Term{acyclic}{$(n-1)$-无上同调}性假设的等式 $R^ig_*(g^*(F))=0$（$0<i\leq n-1$）和 $F\xrightarrow{\sim}g_*g^*(F)$，给出\Term{canonical}{典范}\Term{exact-sequence}{正合列}
\[0\longrightarrow H^n(Y,F)\longrightarrow H^n(X,g^*(F))\longrightarrow H^0(Y,R^ng_*(g^*(F))).\]
\SourcePage{130}
这表明，只需证明 $\xi$ 在第三项中的像为零。此事可逐茎验证，结合第 VIII 讲 5.2，便归结到 $Y$ \Term{strictly-local}{严格局部}的情形。$n=0$、$F$ 为\Term{set}{集合}\Term{sheaf}{层}的情形也同样处理：关于 $g$ 的假设蕴含 $F\to g_*g^*(F)$ \Term{injectivity}{单射}，而 $\xi$ 可视为后一\Term{sheaf}{层}的\Term{section}{截面}；须证明它来自前一\Term{sheaf}{层}的\Term{section}{截面}。这仍可逐茎验证，再应用第 VIII 讲 5.3。同一论证基本上也适用于 $n=1$ 的非交换情形，只需使用第 XII 讲 3.2 的\Term{exact-sequence}{正合列}
\[H^1(Y,F)\longrightarrow H^1(X,g^*(F))\longrightarrow H^0(Y,R^1g_*(g^*(F)))\]
以及第 VIII 讲 5.3。

故以下假设 $Y$ \Term{strictly-local}{严格局部}，从而 $X$ \Term{quasi-compact}{拟紧}且\Term{quasi-separated}{拟分离}。\Term{morphism}{态射} $g:X\to Y$ 满足第 XII 讲 6.10 的条件，因此归结到如下情形：对 $Y$ 的每个异于 $Y$ 的\Term{closed-subset}{闭子集} $Y_1$，$\xi$ 在 $X_1=X\times_Y Y_1$ 上的\Term{restriction}{限制}均属于 $H^n(Y_1,F|_{Y_1})\to H^n(X_1,g^*(F)|_{X_1})$ 的像；而此时须验证，可以找到\Term{finite-morphism}{有限态射} $f:Y'\to Y$，使其在 $Y$ 中的像包含 $Y$ 的一个\Term{nonempty}{非空}\Term{open-subset}{开集}，且 $\xi$ 在 $X'=X\times_Y Y'$ 上的\Term{inverse-image}{逆像} $\xi'$ 属于 $H^n(Y',F')\to H^n(X',g'^*(F'))$ 的像。注意，结合第 VIII 讲 5.5、5.8，上引处的条件 $2^\circ)$ 对此\Term{morphism}{态射} $f:Y'\to Y$ 自动成立。

为找到这样的 $f$，可假设 $Y$ \Term{reduced}{约化}，并取 $Y$ 的一个\Term{maximal-point}{极大点} $y$。设 $\bar y$ 为对应于 $k(y)$ 的一个\Term{separable-closure}{可分闭包}的\Term{geometric-point}{几何点}。则 $\bar y$ 是 $Y$ 上\Term{finite}{有限}整\Term{scheme}{概形} $Y_i$ 的\Term{open-subset}{开集} $U_i$ 的\Term{filtered-projective-limit}{滤过投射极限}，其中每个 $Y_i$ 在 $Y$ 中的像均包含 $y$ 的一个\Term{neighborhood}{邻域}。利用关于 $\xi$ 的假设，可知存在指标 $i$，使 $\xi$ 在 $X\times_Y U_i$ 上的\Term{inverse-image}{逆像}，在 $n\geq1$ 时为零，在 $n=0$ 时属于 $H^0(U_i,F_{U_i})$ 的像。

取 $Y'=Y_i$。注意，可以假设 $U'=U_i$ 是 $y$ 在 $Y$ 中的某个\Term{affine}{仿射}\Term{open-neighborhood}{开邻域} $U$ 的\Term{inverse-image}{逆像}；于是 $\xi$ 在 $X\times_Y Z'$（其中 $Z'=Y'-U'$ 为 $Z=X-U$ 在 $Y'$ 中的\Term{inverse-image}{逆像}）上的\Term{inverse-image}{逆像}属于 $H^n(Z',F_{Z'})$ 的像。将 $Y$ 换成 $Y'$ 后，便归结到如下情形：存在 $Y$ 的\Term{affine}{仿射}\Term{open-subset}{开集} $U$，其补集为 $Z$，使 $\xi$ 在 $X_U$ 及 $X_Z$ 上的\Term{restriction}{限制}分别属于 $H^n(U,F_U)$ 及 $H^n(Z,F_Z)$ 的像；此外，当 $n\geq1$ 时，可假设这两个\Term{restriction}{限制}为零。

当 $n\geq1$ 时，我们将证明在这些条件下 $\xi=0$。为简化记号，把 $X$ 写成 $Y'$，以 $U',Z'$ 表示 $U,Z$ 在 $Y'$ 中的\Term{inverse-image}{逆像}，并以 $F'$ 表示\Term{inverse-image}{逆像} $g^*(F)$。记 $i:U\to Y$、
\SourcePage{131}
$j:Z\to Y$、$i':U'\to Y'$、$j':Z'\to Y'$ 为包含映射，
\[
\begin{tikzcd}
U'\arrow[r,"i'"]\arrow[d,"g_U"']&Y'=X\arrow[r,"j'"]\arrow[d,"g"]&Z'\arrow[d,"g_Z"]\\
U\arrow[r,"i"']&Y\arrow[r,"j"]&Z
\end{tikzcd}
\]
并考虑\Term{sheaf}{层}的\Term{exact-sequence}{正合列}
\[0\longrightarrow F'_{U',Y'}\longrightarrow F'\longrightarrow F'_{Z',Y'}\longrightarrow0.\]
由于 $\xi$ 在 $H^n(Y',F'_{Z',Y'})=H^n(Z',F'_{Z'})$ 中的像为零，$\xi$ 来自 $H^n(Y',F'_{U',Y'})$ 的某个元素；后者在 $U'$ 上的\Term{restriction}{限制}等于 $\xi$ 的\Term{restriction}{限制}，因而为零。$F'_{U',Y'}$ 与 $F_{U,Y}$ 的\Term{inverse-image}{逆像}\Term{isomorphism}{同构}，故归结到形如 $F_{U,Y}$ 的\Term{sheaf}{层}，于是可假设 $F_Z=0$（这已经考虑了 $\xi|_{Z'}=0$ 的假设）。注意，若 $F\to G$ 为阿贝尔 $L$-\Term{torsion-sheaf}{挠层}（相应地，群\Term{sheaf}{层}）的\Term{monomorphism}{单态射}，只需证明 $\xi$ 在 $H^n(Y',G')$ 中的像 $\eta$ 属于 $H^n(Y,G)$ 的像，即为零（第 XII 讲 6.6）。这允许将 $F$ 换成 $i_*(F_U)$，因为 $F_Z=0$ 蕴含\Term{canonical-homomorphism}{典范同态} $F\to i_*(F_U)$ \Term{injectivity}{单射}。由于 $g$ 对 $L$ \Term{locally-acyclic}{局部 $(n-1)$-无上同调}，因 $n\geq1$，它尤其对 $L$ \Term{acyclic}{$0$-无上同调}，从而 $i_*(F_U)$ 在 $Y'$ 上的\Term{inverse-image}{逆像}可与 $i'_*(F'_{U'})$ 等同。故只须证明 $\xi$ 在 $H^n(Y',i'_*(F'_{U'}))$ 中的像 $\eta$ 为零。$\eta$ 在 $U'$ 上的\Term{restriction}{限制}为零，因此只须验证\Term{homomorphism}{同态}
\[H^n(Y',i'_*(F'_{U'}))\longrightarrow H^n(U',F'_{U'})\tag{*}\]
为\Term{monomorphism}{单态射}。在阿贝尔情形，使用 $i'$ 和 $F'_{U'}$ 的 Leray \Term{spectral-sequence}{谱序列}
\[E_2^{p,q}=H^p(Y',R^qi'_*(F'_{U'}))\Longrightarrow H^{p+q}(U',F'_{U'})\]
及关系
\[E_2^{p,q}=0\qquad(0<p\leq n-1,\ q\leq n-1).\tag{**}\]
为验证这些关系，注意 $g$ 对 $L$ \Term{locally-acyclic}{局部 $(n-1)$-无上同调}的假设蕴含\Term{isomorphism}{同构}
\[R^qi'_*(F'_{U'})\simeq g^*(R^qi_*(F_U))\qquad(q\leq n-1).\]
结合 $g$ 对 $L$ \Term{acyclic}{$(n-1)$-无上同调}，即对 $p\leq n-1$ 及 $Y$ 上每个阿贝尔 $L$-\Term{torsion-sheaf}{挠层} $G$ 均有 $H^p(Y,G)\xrightarrow{\sim}H^p(Y',G')$，再用 $p>0$ 时 $H^p(Y,G)=0$，便得 (**)。同样得到
\[E_2^{0,q}\simeq H^0(Y,R^qi_*(F_U))\qquad(q\leq n-1).\tag{***}\]
\SourcePage{132}
所考察的 Leray \Term{spectral-sequence}{谱序列}于是产生一条\Term{exact-sequence}{正合列}，构成下列\Term{commutative-diagram}{交换图}的第二行：
\[
\begin{tikzcd}[column sep=small]
H^0(Y,R^{n-1}i_*(F_U))\arrow[r]\arrow[d,"\sim"']&H^n(Y,i_*(F_U))=0\arrow[d]&\\
H^0(Y',R^{n-1}i'_*(F'_{U'}))\arrow[r]&H^n(Y',i'_*(F'_{U'}))\arrow[r]&H^n(U',F'_{U'}).
\end{tikzcd}
\]
其中第一条竖直\Term{arrow}{箭头}由 (***) 为\Term{bijection}{双射}。由图可读出 (*) \Term{injectivity}{单射}，故此情形的证明完成。在非阿贝尔情形 $n=1$，使用 $n=1$ 时的同一图：
\[
\begin{tikzcd}[column sep=small]
H^0(Y,i_*(F_U))\arrow[r]\arrow[d,"\sim"']&H^1(Y,i_*(F_U))=0\arrow[d]&\\
H^0(Y',i'_*(F'_{U'}))\arrow[r]&H^1(Y',i'_*(F'_{U'}))\arrow[r]&H^1(U',F'_{U'}),
\end{tikzcd}
\]
见第 XII 讲 3.2；由关于 $g$ 的\Term{local-scheme}{局部}及整体\Term{acyclic}{$0$-无上同调}假设，第一条竖直\Term{arrow}{箭头}仍为\Term{bijection}{双射}，故在此情形同样得出 (*) \Term{injectivity}{单射}。

最后处理 $n=0$。注意，根据\SGAsourcegap{此处“根据”之后原拟引用的文献或编号缺失；主源 PDF 第 132 页及补源 PDF 第 192 页均为空白}，给出 $F'$ 的一个\Term{section}{截面} $\xi$，\Term{equivalence}{等价}于给出 $F'_{U'}$ 的\Term{section}{截面} $\xi_{U'}$ 和 $F'_{Z'}$ 的\Term{section}{截面} $\xi_{Z'}$，使 $\xi_{Z'}$ 在自然\Term{homomorphism}{同态}
\[F'_{Z'}\longrightarrow j'^*(i'_*(F'_{U'}))\]
下的像等于 $\xi_{U'}$ 诱导的右边的\Term{section}{截面}。由假设，存在 $F_Z,F_U$ 的\Term{section}{截面} $\eta_Z,\eta_U$，分别诱导 $\xi_{Z'},\xi_{U'}$。只须证明它们满足相应的相容条件，即它们是 $F$ 的同一\Term{section}{截面} $\eta$ 的\Term{restriction}{限制}。须验证 $j^*(i_*(F_U))$ 的两个\Term{section}{截面}相等；由于 $Z'\to Z$ \Term{surjectivity}{满射}（$g$ 按假设\Term{acyclic}{$(-1)$-无上同调}，因而\Term{surjectivity}{满射}），只需验证 $g_Z^*(j^*(i_*(F_U)))=j'^*(g^*(i_*(F_U)))$ 的相应两个\Term{section}{截面}相等。通过\Term{base-change}{换基}\Term{homomorphism}{同态}
\[g^*(i_*(F_U))\longrightarrow i'_*(g_U^*(F_U))=i'_*(F'_{U'}),\tag{*}\]
所考虑的\Term{sheaf}{层}映入 $j'^*(i'_*(F'_{U'}))$；后一\Term{homomorphism}{同态}\Term{injectivity}{单射}，因为由 $g$ \Term{locally-acyclic}{局部 $(-1)$-无上同调}的假设，(*) \Term{injectivity}{单射}。所以只须验证 $j'^*(i'_*(F'_{U'}))$ 中这两个\Term{section}{截面}的像相等。立即可见，这正是前面所考察的两个\Term{section}{截面}，而它们按假设相等。这就完成了证明。
\end{proof}
\textbf{注记 1.18。}
可以给出 1.15 的一个加强，从而在 $Y$ \Term{locally-noetherian}{局部 Noether}时，为 1.15 提供明显更简单的证明。
\SourcePage{133}
为此，除了关于 $g,F$ 的预备假设外，再假设存在\Term{finite}{有限}多个 $Y$ 的\Term{geometric-point}{几何点} $f_i:\bar y_i\to Y$，使\Term{morphism}{态射}
\[F\longrightarrow\prod_i f_{i*}(f_i^*(F))\]
\Term{injectivity}{单射}（当 $Y$ 为 Noether 且 $F$ \Term{constructible}{可构造}时，此假设自动满足，第 IX 讲 2.14）。我们断言，在这些条件下，判据 1.15 只取这些\Term{geometric-point}{几何点} $\bar y_i$ 便已成立。事实上，由第 XII 讲 6.6，可假设 $F$ 本身形如 $f_*(M_{\bar y})$，其中 $M$ 是\Term{set}{集合}，相应地是群，或阿贝尔 $L$-\Term{torsion-group}{挠群}，而 $f:\bar y\to Y$ 是\Term{geometric-point}{几何点}。设 $f':X'=X_{\bar y}\to X$ 为\Term{canonical}{典范}\Term{morphism}{态射}。为确定起见，假设 $n\geq1$；关于 $g$ 的\Term{local-scheme}{局部}假设蕴含
\[g^*(f_*(M_{\bar y}))\simeq f'_*(M_{X'}),\qquad R^if'_*(M_{X'})=0\quad(1\leq i\leq n-1),\]
第二个关系来自其左端\Term{isomorphism}{同构}于 $g^*(R^if_*(M_{\bar y}))$；后者为零，因为显然 $i\geq1$ 时 $R^if_*(M_{\bar y})=0$。利用 $f'$ 的 Leray \Term{spectral-sequence}{谱序列}，由这些关系推出
\[H^n(X,g^*(F))\simeq H^n(X,f'_*(M_{X'}))\longrightarrow H^n(X',M_{X'})\]
\Term{injectivity}{单射}，因而确实证明了：$H^n(X,F)$ 中诱导 $H^n(X_{\bar y},F)$ 的零元素的元素 $\xi$ 为零。为由此在 $Y$ \Term{locally-noetherian}{局部 Noether}时给出 1.15 的简化证明，注意可归结到 $Y$ 为 Noether 的情形；再利用第 IX 讲 2.7.2 和第 VII 讲 5 作一次容易的取\Term{limit}{极限}，便归结到 $F$ 还\Term{constructible}{可构造}的情形，此时适用刚才证明的断言。

\SGAsection{2}{光滑态射的局部无上同调性}
定理如下。
\begin{SGAresult}{定理}{2.1}
设 $g:X\to Y$ 为\Term{smooth}{光滑}\Term{morphism}{态射}，$L\subset\mathbb P$ 为 $X$ 的\Term{residue-characteristic}{剩余特征}集的补集。则 $g$ 对 $L$ \Term{locally-acyclic}{局部无上同调}且对 $L$ \Term{locally-one-aspherical}{局部 $1$-非球面}。
\end{SGAresult}
先记下一个推论。
\begin{SGAresult}{推论}{2.2}
设 $g:\mathbb E_Y^m\to Y$ 为\Term{affine-space}{仿射空间}的\Term{structural-morphism}{结构态射}。则 $g$ 对 $Y$ 的\Term{residue-characteristic}{剩余特征}集的补集 $L$ \Term{acyclic}{无上同调}且\Term{one-aspherical}{$1$-非球面}。
\end{SGAresult}
事实上，对 $m$ 归纳并用 1.9 (i)，归结到 $m=1$。由 2.1，$g$ 对 $L$ \Term{locally-acyclic}{局部 $n$-无上同调}且对 $L$ \Term{locally-one-aspherical}{局部 $1$-非球面}，所以 1.16 的\Term{local-scheme}{局部}条件满足。因此只须证明 $\mathbb E_Y^1$ 的\Term{geometric-fiber}{几何纤维}对 $L$ \Term{acyclic}{无上同调}且\Term{one-aspherical}{$1$-非球面}，即当 $k$ 为\Term{characteristic}{特征} $p\notin L$ 的\Term{separably-closed}{可分闭}域时，$\operatorname{Spec}k[t]$
\SourcePage{134}
对 $L$ \Term{acyclic}{无上同调}且\Term{one-aspherical}{$1$-非球面}。\Term{acyclic}{无上同调}性由第 IX 讲 4.6 得出，\Term{one-aspherical}{$1$-非球面}性则是熟知的。更一般地，设 $X'$ 为 $X=\mathbb P_k^1$ 的\Term{normal}{正规}\Term{irreducible}{不可约} Galois \Term{covering}{覆盖}，它在点 $\infty$ 之外\Term{unramified}{非分歧}，并在该点\Term{tamely-ramified}{温分歧}。证明此\Term{covering}{覆盖}的次数 $n$ 等于 $1$。有 Hurwitz 公式
\[2(g'-1)=2n(g-1)+\deg\mathfrak d_{X'/X},\]
其中 $g,g'$ 分别为 $X,X'$ 的\Term{genus}{亏格}，$\mathfrak d_{X'/X}$ 为不同\Term{divisor}{除子}。有 $g=0$，而\Term{tamely-ramified}{温分歧}假设蕴含 $\deg\mathfrak d_{X'/X}\leq n-1$，故 $2(g'-1)\leq-(n+1)$，从而 $n\leq2(1-g')-1$；由于 $g'\geq0$，得 $n\leq1$，所以 $n=1$。证毕。
\begin{proof}[2.1 的证明]
该断言对\Term{etale-topology}{平展拓扑}在 $X,Y$ 上均为\Term{local-scheme}{局部}的，故可假设（SGA 1 第 II 讲 1.1）$g$ 为\Term{affine-space}{仿射空间} $\mathbb E_Y^m$ 的\Term{structural-morphism}{结构态射}，且 $Y$ \Term{affine}{仿射}。应用\Term{locally-acyclic}{局部无上同调}性的\Term{transitivity}{传递性} 1.13 (i)，以及 $\mathbb E_Y^m=\mathbb E_{\mathbb E_Y^1}^{m-1}$，立即归结到 $m=1$，即 $X=\operatorname{Spec}\mathscr O_Y[t]$ 的情形。

应用 1.18。只须证明：由一对\Term{geometric-point}{几何点} $\bar x,\bar y$ 得到的\Term{strictly-local}{严格局部}\Term{scheme}{概形}\Term{morphism}{态射} $\widetilde g:\widetilde X\to\widetilde Y$，其中 $x$ 在\Term{fiber}{纤维} $X_y$ 中为\Term{closed-point}{闭点}，其\Term{geometric-fiber}{几何纤维}均对 $L$ \Term{acyclic}{无上同调}且对 $L$ \Term{one-aspherical}{$1$-非球面}。显然可假设 $Y=\widetilde Y$。设 $Y'\to Y$ 为\Term{finite}{有限}、\Term{local-scheme}{局部}、\Term{surjectivity}{满射}的\Term{morphism}{态射}，于是 $Y'$ \Term{strictly-local}{严格局部}；设 $y'$ 为 $Y'$ 中 $y$ 上方的点，$x'$ 为 $X'$ 中 $x$ 及 $y'$ 上方的点。由于 $X'\to X$ \Term{finite}{有限}，$X'$ 在 $x'$ 处的\Term{strict-localization}{严格局部化}正是 $\widetilde X'=\widetilde X\times_Y Y'$。因此，$\widetilde X'/Y'$ 的每个\Term{geometric-fiber}{几何纤维}均由 $\widetilde X/Y$ 的一个\Term{geometric-fiber}{几何纤维}经过底域的\Term{algebra}{代数}\Term{group-extension}{扩张}（必为\Term{radicial}{根性}\Term{group-extension}{扩张}）得到；此外，$\widetilde X/Y$ 的每个\Term{geometric-fiber}{几何纤维}上方至少有 $\widetilde X'/Y'$ 的一个\Term{geometric-fiber}{几何纤维}。故由第 VIII 讲 1.1，可将 $Y$ 换成 $Y'$；适当选择 $Y'$，便归结到 $k(x)=k(y)$ 的情形。此时 $x$ 为\Term{fiber}{纤维} $X_y=\operatorname{Spec}k(y)[t]$ 的\Term{rational-point}{有理点}，例如 $t=a_0$。将 $a_0$ \Term{lifting}{提升}为 $\mathscr O_{\widetilde Y}$ 的\Term{section}{截面} $a$，再作“\Term{translation}{平移}”$t_1=t-a$，便归结到 $x$ 为 $\operatorname{Spec}k(y)[t]$ 的点 $\{t=0\}$ 的情形。令 $A=\Gamma(\widetilde Y,\mathscr O_{\widetilde Y})$，于是 $A$ 为\Term{strictly-local}{严格局部}环。
\end{proof}
\begin{SGAresult}{记号}{2.3}
设 $A$ 为\Term{local-scheme}{局部}\Term{henselian}{Hensel 环}。以 $A\{t\}$ 表示环 $A[t]$ 在由 $\operatorname{rad}A$ 及 $t$ 生成的\Term{prime-ideal}{素理想}处的\Term{henselization}{Hensel 化}。
\end{SGAresult}
\begin{SGAresult}{推论}{2.4}
设 $A$ 为\Term{henselian}{Hensel 环}，$A'$ 为\Term{finite}{有限} $A$-\Term{algebra}{代数}。则
\[A'\{t\}\simeq A'\otimes_A A\{t\}.\]
\end{SGAresult}
这立即由第 VIII 讲 4.1 得出。注意，若 $A$ \Term{strictly-local}{严格局部}，则 $A\{t\}$ 也如此。

定理的证明归结到下列引理。
\SourcePage{135}
\begin{SGAresult}{引理}{2.5}
设 $A$ 为\Term{residue-characteristic}{剩余特征} $p$ 的\Term{henselian}{Hensel 环}。则 $\operatorname{Spec}A\{t\}$ 在 $\operatorname{Spec}A$ 上的每个\Term{geometric-fiber}{几何纤维}，均对 $L=\mathbb P-\{p\}$ \Term{acyclic}{无上同调}且\Term{one-aspherical}{$1$-非球面}。
\end{SGAresult}
注意，由于 $A\{t\}$ 是 $A[t]$ 的\Term{henselization}{Hensel 化}，\Term{fiber}{纤维} $\operatorname{Spec}(A\{t\}\otimes_A K)$（$K$ 为 $A$ 的一个\Term{residue-field}{剩余域}）是“直线”$\operatorname{Spec}K[t]$ 上\Term{affine}{仿射}\Term{finite-type}{有限型}\Term{scheme}{概形} $X_i$ 的\Term{inverse-limit}{投射极限}，即\Term{affine}{仿射}\Term{algebraic-curve}{代数曲线}的\Term{limit}{极限}。所以 $\operatorname{Spec}A\{t\}$ 的每个\Term{geometric-fiber}{几何纤维}都是\Term{separably-closed}{可分闭}域 $\overline K$ 上\Term{affine}{仿射}曲线 $(X_i)_{\overline K}$ 的\Term{limit}{极限}；由第 IX 讲 5.7 及第 VII 讲 5.7，这些\Term{geometric-fiber}{几何纤维}的\Term{cohomological-dimension}{上同调维数}均 $\leq1$。故只须验证它们对 $L$ \Term{one-aspherical}{$1$-非球面}，即这样的\Term{geometric-fiber}{几何纤维} $Z$ \Term{connected}{连通}、\Term{nonempty}{非空}，且对每个\Term{finite}{有限} $L$-群 $G$，$H^1(Z,G)=0$。

写成 $A=\varinjlim B_\alpha$，其中 $B_\alpha$ 为\Term{finite-type}{有限型} $\mathbb Z$-\Term{algebra}{代数}。\Term{morphism}{态射} $B_\alpha\to A$ 通过 $B_\alpha$ 在 $\operatorname{rad}A$ 诱导的\Term{prime-ideal}{素理想}处的\Term{henselization}{Hensel 化} $A_\alpha$ \Term{resolution}{分解}，故 $A=\varinjlim A_\alpha$。容易看出，同样有 $A\{t\}=\varinjlim A_\alpha\{t\}$。再设 $\bar y$ 为 $\operatorname{Spec}A$ 的\Term{geometric-point}{几何点}，$X_{\bar y}$ 为 $S=\operatorname{Spec}A\{t\}$ 在 $\operatorname{Spec}A$ 上 $\bar y$ 处的\Term{geometric-fiber}{几何纤维}；$\bar y_\alpha$ 为 $\bar y$ 诱导的 $\operatorname{Spec}A_\alpha$ 的\Term{geometric-point}{几何点}，$X_{\alpha\bar y_\alpha}$ 为 $\operatorname{Spec}A_\alpha\{t\}$ 在 $\bar y_\alpha$ 处的\Term{geometric-fiber}{几何纤维}。有 $X_{\bar y}=\varprojlim X_{\alpha\bar y_\alpha}$。所以，证明 2.5 时，由 1.7 (ii)（应用于\Term{morphism}{态射} $X_{\alpha\bar y_\alpha}\to\bar y_\alpha$）可归结到 $A=A_\alpha$，即 $A$ 为某个\Term{finite-type}{有限型} $\mathbb Z$-\Term{algebra}{代数}在给定\Term{prime-ideal}{素理想}处的\Term{henselization}{Hensel 化}的情形。以下作此假设。回顾这样的\Term{algebra}{代数}是\Term{excellent}{优良}环（EGA IV 7.8.3 (iii)、7.8.6 (i)），故 $A_\alpha$ 也\Term{excellent}{优良}（由定义容易看出，参见 EGA IV 18.7.6）。

验证\Term{acyclic}{$0$-无上同调}性。（更一般的结果见附录。）由于 $A\{t\}/A$ 有\Term{section}{截面}“$t=0$”，\Term{geometric-fiber}{几何纤维} $Z$ \Term{nonempty}{非空}。为证明它\Term{connected}{连通}，由 2.4 把 $A$ 换成一个商，归结到 $A$ 整且 $Z$ 为 $\operatorname{Spec}A$ 的\Term{generic-point}{泛点}上方的\Term{geometric-fiber}{几何纤维}的情形。只须证明，对 $A$ 的\Term{field-of-fractions}{分式域} $K$ 的任意\Term{finite-separable-extension}{有限可分扩张} $K'$，\Term{scheme}{概形} $Z_{K'}$（$Z=\operatorname{Spec}(A\{t\}\otimes_A K)$）\Term{connected}{连通}。设 $A'$ 为 $A$ 在 $K'$ 中的\Term{normalization}{正规化}（它是\Term{finite}{有限} $A$-\Term{algebra}{代数}，EGA IV 7.8.2），$Z'$ 为 $\operatorname{Spec}A'\{t\}$ 在 $A'$ 上的\Term{generic-fiber}{泛纤维}。由 2.4，有 $A'\{t\}\simeq A'\otimes_A A\{t\}$，故 $Z_{K'}\simeq Z'$。将 $A$ 换成 $A'$。于是只须证明 $Z$ 整，而这又只须证明 $A\{t\}$ 整。由于 $A$ \Term{normal}{正规}，$A[t]$ 也\Term{normal}{正规}，故 $A\{t\}$ 也\Term{normal}{正规}（SGA 1 第 I 讲 9.5 (i)），因而得到结论。

因此，2.5 的证明，从而 2.1 的证明，归结到下述引理，其证明将在下一节给出。
\begin{SGAresult}{引理}{2.6}
设 $A$ 为\Term{excellent}{优良}\Term{henselian}{Hensel 环}，$\overline Z$ 为 $\operatorname{Spec}A\{t\}/\operatorname{Spec}A$ 的一个\Term{geometric-fiber}{几何纤维}。则 $\overline Z$ 的每个主 Galois \Term{covering}{覆盖} $\overline Z'$，若其群 $G$ 的阶数与 $A$ 的\Term{residue-characteristic}{剩余特征} $p$ 互素，便是\Term{trivial-covering}{平凡覆盖}。
\end{SGAresult}

\SourcePage{136}
\SGAsection{3}{主要引理的证明}
将 $A$ 换成\Term{finite}{有限}\Term{algebra}{代数} $A'$（参见 2.4），把 2.6 归结到 $A$ 整且 $\overline Z$ 为 $\operatorname{Spec}A\{t\}$ 在 $\operatorname{Spec}A$ 上的几何\Term{generic-fiber}{泛纤维}的情形。设 $K$ 为 $A$ 的\Term{field-of-fractions}{分式域}，$Z=\operatorname{Spec}(A\{t\}\otimes_A K)$。若 $\overline Z'/\overline Z$ 为主 Galois \Term{covering}{覆盖}，则可将其\Term{descent}{下降}到某个 $Z_{K'}$，其中 $K'$ 是 $K$ 的适当\Term{finite-separable-extension}{有限可分扩张}。根据 2.4，将 $A$ 换成它在 $K'$ 中的\Term{normalization}{正规化}，这是\Term{finite}{有限} $A$-\Term{algebra}{代数}（EGA IV 7.8.3 (vi)），便归结到证明引理的下述\Term{form}{形式}。
\begin{SGAresult}{引理}{3.1}
沿用 2.6 的记号，假设 $A$ \Term{normal}{正规}，$Z'/Z$ 为\Term{etale}{平展}主 Galois \Term{covering}{覆盖}，其群 $G$ 的阶数与 $A$ 的\Term{residue-characteristic}{剩余特征} $p$ 互素。则 $Z'$ 由 $K$ 的一个 Galois \Term{group-extension}{扩张} $K'$ 诱导，即（SGA 1 第 IX 讲 6.2）由 $Z'$ 得到的 $\overline Z$ 的\Term{covering}{覆盖} $\overline Z'$ 平凡。
\end{SGAresult}
\begin{proof}
设 $B'$ 为 $A\{t\}$ 在 $Z'$ 的\Term{rational-function-ring}{有理函数环}中的\Term{normalization}{正规化}，考虑 $A\{t\}$ 中使\Term{algebra}{代数} $B'/A\{t\}$ 在其处有\Term{ramified}{分歧}的\Term{height}{高度} $1$ \Term{prime-ideal}{素理想} $P$。由于 $Z'/Z$ \Term{etale}{平展}，$P$ 不在 $\operatorname{Spec}A$ 的\Term{generic-point}{泛点}上方，即 $\mathfrak p=P\cap A\ne0$。故 $A\{t\}/\mathfrak p A\{t\}=(A/\mathfrak p)\{t\}$ 整，且\Term{dimension}{维数}至多为 $\dim A$。$A\{t\}$ \Term{excellent}{优良}（EGA IV 7.8.6 (i)、18.7.6），因而\Term{catenary}{链状}；故
\[\dim\operatorname{Spec}(A\{t\}/P)=(\dim A+1)-1=\dim A\geq\dim(A\{t\}/\mathfrak p A\{t\}),\]
由于 $P\supset\mathfrak p A\{t\}$，遂有 $P=\mathfrak p A\{t\}$。

于是，设 $\mathfrak p_1,\ldots,\mathfrak p_m$ 为使 $B'/A\{t\}$ 在\Term{ideal}{理想} $\mathfrak p_iA\{t\}$ 处有\Term{ramified}{分歧}的 $A$ 的\Term{height}{高度} $1$ \Term{prime-ideal}{素理想}；设 $s$ 为 $B'$ 在各 $\mathfrak p_iA\{t\}$ 上方\Term{prime-ideal}{素理想}处的\Term{inertia-group}{惯性群} $G_i\subset G$（SGA 第 IV 讲 2）之阶数的最大公因数。它是整除 $G$ 的阶数的整数，因而与 $p$ 互素。设 $f\in A$ 为在每个 $\mathfrak p_i$ 处恰有一阶零点的函数（这样的 $f$ 存在，例如可由 Bourbaki，\textit{Alg. Comm.}，第 II 章，\S3，命题 17 的推论得出）。根据 2.4，将 $A$ 换成它在 $K$ 的\Term{finite-extension}{有限扩张} $K_1=K[x]/(x^s-f)$ 中的\Term{normalization}{正规化} $A_1$，并将 $B'$ 换成 $A_1\otimes_A B'$ 的\Term{normalization}{正规化}，由 Abhyankar 引理（SGA 1 第 X 讲 3.6），便归结到 $B'$ 在 $A\{t\}$ 的任何\Term{height}{高度} $1$ \Term{prime-ideal}{素理想}处都无\Term{ramified}{分歧}的情形。

于是\Term{algebra}{代数} $B'/A\{t\}$ 在 $\operatorname{Spec}A$ 每个\Term{codimension}{余维} $1$ 点上方的 $\operatorname{Spec}A\{t\}$ 的\Term{fiber}{纤维}上也\Term{etale}{平展}。事实上，设 $\mathfrak p$ 为 $A$ 的\Term{height}{高度} $1$ \Term{prime-ideal}{素理想}。由于 $A$ \Term{normal}{正规}，$A_{\mathfrak p}$ 为\Term{discrete-valuation-ring}{离散赋值环}，故\Term{regular}{正则}；所以 $A\{t\}\otimes_A A_{\mathfrak p}$ 也\Term{regular}{正则}，因为它是在 $A_{\mathfrak p}[t]$ 上\Term{etale}{平展}的\Term{scheme}{概形}的\Term{limit}{极限}。由假设，\Term{algebra}{代数} $B'/A\{t\}$ 在 $\operatorname{Spec}(A\{t\}\otimes_A A_{\mathfrak p})$ 的每个\Term{codimension}{余维} $\leq1$ 的点处\Term{etale}{平展}，故由 Zariski--Nagata \Term{purity-theorem}{纯性定理}（SGA 1 第 X 讲 3.1；证明例如见 SGA 2 第 X 讲 3.4），它处处\Term{etale}{平展}。

根据 2.4，将 $A$ 换成适当的\Term{finite}{有限} $A$-\Term{algebra}{代数}，还可归结到\Term{algebra}{代数} $B'/tB'$ 在 $A\{t\}/tA\{t\}\simeq A$ 上，于 $\operatorname{Spec}A$ 的\Term{generic-point}{泛点}上方\Term{completely-split}{完全分裂}的情形。所以只须证明，在这些条件下，$A\{t\}$ 上的\Term{algebra}{代数} $B'$ “平凡”。由于 $A\{t\}$ 为\Term{henselian}{Hensel 环}，只须证明
\SourcePage{137}
$B'$ 在 $A\{t\}$ 上\Term{etale}{平展}；事实上，\Term{algebra}{代数} $B'/tB'$ 在 $\operatorname{Spec}A$ 的\Term{generic-point}{泛点}上方\Term{splitting}{分裂}，蕴含它在 $\operatorname{Spec}A$ 上\Term{splitting}{分裂}（SGA 1 第 I 讲 10.1），从而 $B'$ 在 $A\{t\}$ 的\Term{maximal-ideal}{极大理想}上方的\Term{residue-extension}{剩余扩张}平凡，再应用 EGA IV 18.5.14。因此归结到下述引理。
\end{proof}
\begin{SGAresult}{引理}{3.2}
设 $B$ 为\Term{excellent}{优良}\Term{normal}{正规}\Term{local-ring}{局部环}，$0\ne t\in\operatorname{rad}B$，且 $B/tB$ \Term{normal}{正规}。设 $B'$ 为包含 $B$ 的\Term{finite}{有限}整\Term{normal}{正规} $B$-\Term{algebra}{代数}，满足：
\begin{enumerate}
\item[(i)] $B'/tB'$ 作为 $B/tB$-\Term{algebra}{代数}，在 $\operatorname{Spec}(B/tB)$ 的\Term{maximal-point}{极大点}处\Term{completely-split}{完全分裂}。
\item[(ii)] $B'$ 在 $B$ 上的\Term{algebra}{代数}，在 $\operatorname{Spec}(B/tB)$ 的\Term{codimension}{余维} $1$ 点处\Term{etale}{平展}。
\end{enumerate}
则 $B'$ 在 $B$ 上\Term{etale}{平展}。
\end{SGAresult}
\begin{proof}
由 (ii)，当 $\dim B\leq2$，即 $\dim B/tB\leq1$ 时，断言显然。假设 $B'/B$ 不\Term{etale}{平展}。设 $x$ 为 $\operatorname{Spec}(B/tB)$ 中使 $B'/B$ 在其上方不\Term{etale}{平展}的点集的一个\Term{maximal-point}{极大点}，$C$ 为 $\operatorname{Spec}B$ 在 $x$ 处的\Term{local-ring}{局部环}。则 $C$ 仍为\Term{excellent}{优良}\Term{normal}{正规}环（EGA IV 7.8.2），且 $0\ne t\in\operatorname{rad}C$，而 (i)、(ii) 对 $C$-\Term{algebra}{代数} $C'=C\otimes_B B'$ 成立。因此可假设 $B=C$，即 $B'$ 在 $\operatorname{Spec}(B/tB)$ 的每点上方都在 $B$ 上\Term{etale}{平展}，唯一可能的例外是原点。由于 $B/tB$ \Term{normal}{正规}（其实\Term{unibranch}{单支}就够了），条件 (i) 蕴含\Term{algebra}{代数} $B'/tB'$ 在原点之外\Term{completely-split}{完全分裂}。此外可假设 $\dim B\geq3$。故归结到下述引理。
\end{proof}
\begin{SGAresult}{引理}{3.3}
设 $B$ 为\Term{excellent}{优良}\Term{normal}{正规}\Term{local-ring}{局部环}，$0\ne t\in\operatorname{rad}B$，$\dim B\geq3$。设 $B'$ 为包含 $B$ 的\Term{finite}{有限}\Term{normal}{正规}整 $B$-\Term{algebra}{代数}，且 $B'/tB'$ 作为 $B/tB$-\Term{algebra}{代数}在原点之外\Term{completely-split}{完全分裂}。则 $B'$ 在 $B$ 上\Term{etale}{平展}。
\end{SGAresult}
\begin{proof}
设 $X=\operatorname{Spec}B$，$X_1=X-\{x_0\}$，其中 $x_0$ 是\Term{closed-point}{闭点}；设 $i:U\to X$ 为使\Term{algebra}{代数} $B'/B$ 在其处\Term{etale}{平展}的 $x\in X$ 构成的\Term{open-subset}{开集}。则 $(X-U)\cap V(t)=\{x_0\}$ 的\Term{dimension}{维数}为 $0$。$V(t)$ 由一个方程定义，故 $\dim(X-U)\leq1$（EGA $0_{\mathrm{IV}}$ 16.2.5）。由于 $X$ \Term{catenary}{链状}，$x\in X-U$ 蕴含 $\dim\mathscr O_{X,x}\geq\dim B-1\geq2$，即 $\operatorname{codim}(X-U,X)\geq2$。

设 $X'=\operatorname{Spec}(B')$，$U'$ 为 $U$ 在 $X'$ 中的\Term{inverse-image}{逆像}。$X$ \Term{universally-catenary}{普适链状}，故 $\operatorname{codim}(X-U,X)\geq2$ 蕴含 $\operatorname{codim}(X'-U',X')\geq2$（EGA IV 5.6.10）；又 $X'$ \Term{normal}{正规}，所以
\[B'=\Gamma(X',\mathscr O_{X'})\xrightarrow{\sim}\Gamma(U',\mathscr O_{U'}).\]
设 $\mathbf B'=\widetilde{B'}$ 为在 $X$ 上定义 $X'$ 的 $\mathscr O_X$-\Term{algebra}{代数}\Term{sheaf}{层}，$C$ 为它在 $U$ 上的\Term{restriction}{限制}。则上式\Term{equivalence}{等价}于
\[\mathbf B'\xrightarrow{\sim}i_*(C).\]
\SourcePage{138}
因此 3.3 归结到下述断言（其中已不再出现 $B'$）：设 $B,t$ 如 3.3，$X=\operatorname{Spec}(B)$，$Y=\operatorname{Spec}(B/tB)$，$U$ 为 $Y_1=Y-\{x_0\}$ 在 $X_1=X-\{x_0\}$ 中的一个\Term{open-neighborhood}{开邻域}，$U'$ 为 $U$ 的\Term{etale}{平展}\Term{covering}{覆盖}，其在 $Y_1$ 上的\Term{restriction}{限制}\Term{completely-split}{完全分裂}，$C$ 为定义 $U'$ 的 $\mathscr O_{U'}$-\Term{algebra}{代数}\Term{sheaf}{层}。则以 $i:U\to X$ 表示\Term{canonical}{典范}\Term{immersion}{浸入}，$i_*(C)$ 是 $X$ 上的\Term{coherent}{凝聚}\Term{etale}{平展}\Term{algebra}{代数}（或者，\Term{equivalence}{等价}地，由前述 $x\in X-U$ 时 $\operatorname{depth}\mathscr O_{X,x}\geq2$ 及 EGA IV 5.10.5，$U'$ 可延拓为 $X$ 的\Term{etale}{平展}\Term{covering}{覆盖} $X'$）。

设 $\widehat X=\operatorname{Spec}\widehat B$，其中 $\widehat B$ 为 $B$ 的\Term{completion}{完备化}。则 $\widehat B$ \Term{normal}{正规}（EGA IV 7.8.3 (v)），且 $\widehat B/B$ \Term{faithfully-flat}{忠实平坦}。设
\[
\begin{tikzcd}
\widehat U\arrow[r,"i"]\arrow[d,"f_1"']&\widehat X\arrow[d,"f"]\\
U\arrow[r,"i"']&X
\end{tikzcd}
\]
为由 $\widehat X\to X$ 得到的\Term{cartesian}{笛卡尔}图。由于 $f$ \Term{faithfully-flat}{忠实平坦}，\Term{functor}{函子} $i_*$ 与\Term{base-change}{换基} $f$ 交换（EGA IV 2.3.1）；由 EGA IV 17.7.1，只须证明 $i_*(\widehat C)$ 是 $\widehat X$ 上的\Term{etale}{平展}\Term{algebra}{代数}。所以可将 $B$ 换成 $\widehat B$，即假设 $B$ \Term{complete}{完备}。此时我们将证明 $U'$ 为 $U$ 的\Term{completely-split}{完全分裂}\Term{covering}{覆盖}（这显然蕴含它可延拓为 $X$ 的\Term{etale}{平展}\Term{covering}{覆盖}，从而完成证明）。现在此断言是“\Term{local-scheme}{局部} Lefschetz 理论”（SGA 2 第 X 讲 2.1 (i)、2.5）的直接推论；再次考虑当 $\dim\overline{\{x\}}=1$ 时 $\operatorname{depth}\mathscr O_{X,x}\geq2$，便可立即看出。
\end{proof}
\textbf{注记 3.4。}
基本上沿用同一证明，可以在更一般的条件下证明 3.2、3.3：无需假设 $B$ \Term{excellent}{优良}且\Term{normal}{正规}，只须假设\Term{completion}{完备化} $\widehat B$ \Term{normal}{正规}（由 EGA IV 7 的结果，此条件在\Term{localization}{局部化}下\Term{stable}{稳定}），甚至只须 $B$ 对 $(tB)$-\Term{adic-topology}{进拓扑}的\Term{completion}{完备化}\Term{normal}{正规}且为\Term{regular}{正则}环的商。

\SGAsection{4}{附录：局部 0-无上同调性的一个判据}
本节的结果在本讨论班后续部分不再使用。
\begin{SGAresult}{定理}{4.1}
设 $g:X\to Y$ 为\Term{scheme}{概形}的\Term{morphism}{态射}。假设 $g$ \Term{flat}{平坦}且有\Term{separable}{可分}\Term{fiber}{纤维}，并且 $X,Y$ \Term{locally-noetherian}{局部 Noether}；相应地，假设 $g$ \Term{locally-finitely-presented}{局部有限呈示}。则 $g$ \Term{universally-locally-acyclic}{普适局部 $0$-无上同调}。
\end{SGAresult}
\begin{proof}
显然可假设 $X,Y$ \Term{affine}{仿射}；在\Term{lie-bracket}{括号}中的相应情形，利用 EGA IV 8 的标准方法，并在此用 EGA IV 9.7.7，可归结到 $Y$ 从而 $X$ 为 Noether 的情形。所以只须证明第一个陈述。
\SourcePage{139}
由于假设在\Term{finite-type}{有限型}\Term{base-change}{换基}下\Term{stable}{稳定}，结合取\Term{limit}{极限} 1.13 (ii)，只须证明 $g$ \Term{locally-acyclic}{局部 $0$-无上同调}；由 1.17，这又归结到证明：若进一步假设 $g$ \Term{strictly-local}{严格局部}，则 $g$ 的\Term{geometric-fiber}{几何纤维}均\Term{connected}{连通}。而这正是 EGA IV 18.9.8 的内容。
\end{proof}
\begin{SGAresult}{推论}{4.2}
设 $g:X\to Y$ 为\Term{scheme}{概形}的\Term{morphism}{态射}，$Y$ \Term{discrete}{离散}。则 $g$ \Term{universally-locally-acyclic}{普适局部 $0$-无上同调}。
\end{SGAresult}
\begin{proof}
可以假设 $Y$ 为某个\Term{perfect-field}{完美域} $k$ 的谱（第 VIII 讲 1.1）。另一方面，可假设 $X$ \Term{affine}{仿射}，从而为 $k$ 上\Term{finite-type}{有限型}\Term{affine}{仿射}\Term{scheme}{概形}的\Term{inverse-limit}{投射极限}；结合 1.13 (ii)，归结到 $X$ 在 $k$ 上\Term{finite-type}{有限型}的情形。最后，显然可假设 $X$ \Term{reduced}{约化}；由于 $k$ \Term{perfect}{完美}，$X$ 在 $k$ 上\Term{separable}{可分}。结论于是由 4.1 得出。
\end{proof}
\begin{SGAresult}{推论}{4.3}
假设 $g$ \Term{surjectivity}{满射}，并满足 4.1 或 4.2 的条件；在 4.1 非\Term{lie-bracket}{括号}中的情形，还假设 $g$ 或者\Term{quasi-compact}{拟紧}且\Term{quasi-separated}{拟分离}，或者\Term{universally-open}{普适开}。则对于可变\Term{prescheme}{预概形}上的\Term{etale}{平展}\Term{sheaf}{层}的\Term{fibred-category}{纤维范畴}，$g$ 是\Term{universally-effective-descent}{普适有效下降}\Term{morphism}{态射}。
\end{SGAresult}
\begin{proof}
在所考虑的情形，$g$ \Term{universally-submersive}{普适淹没}；由第 VIII 讲 9.1，$g$ 是所考察\Term{fibred-category}{纤维范畴}的\Term{universal-descent}{普适下降}\Term{morphism}{态射}。这使\Term{effective}{有效}性问题可归结到 $Y$ \Term{affine}{仿射}的情形。此外，在各情形中可见，$X$ 可由\Term{finite}{有限}多个\Term{affine}{仿射}\Term{open-subset}{开集} $X_i$ \Term{covering}{覆盖}，使其像\Term{covering}{覆盖} $Y$。将 $X$ 换成各 $X_i$ 的\Term{disjoint-union}{不交并}，便可假设 $X$ \Term{affine}{仿射}，从而尤其在 $Y$ 上\Term{quasi-compact}{拟紧}且\Term{quasi-separated}{拟分离}。最后利用第 VIII 讲 9.4.1 的论证，并使用第 XVI 讲 1.1 的事实即可：对于\Term{quasi-compact}{拟紧}且\Term{quasi-separated}{拟分离}的\Term{morphism}{态射} $g$，$g_*(F)$（$F$ 为 $X$ 上的\Term{etale}{平展}\Term{sheaf}{层}）的形成与任意\Term{locally-acyclic}{局部 $0$-无上同调}的\Term{base-change}{换基} $Y'\to Y$ 交换。
\end{proof}
\textbf{注记 4.4。}
a) 由此得到第 VIII 讲 9.4 d) 中曾悬而未决的证明，作为 4.3 的特例。另有一个更容易且不用 EGA IV 18.9.8 这一较微妙结果的证明：注意，由第 VIII 讲 9.1，只须考虑由域\Term{group-extension}{扩张} $K/k$ 诱导的\Term{morphism}{态射} $\operatorname{Spec}(K)\to\operatorname{Spec}(k)$；而 $\operatorname{Spec}(K)$ 在 $\operatorname{Spec}(k)$ 上\Term{universally-locally-acyclic}{普适局部 $0$-无上同调}，因为如第 XVI 讲 1.5 将见，它甚至对 $L=\mathbb P-\{p\}$（$p=\operatorname{char}k$）\Term{universally-acyclic}{普适无上同调}。

b) 我们不知道，每个\Term{quasi-compact}{拟紧}且\Term{quasi-separated}{拟分离}、\Term{surjectivity}{满射}并\Term{locally-acyclic}{局部 $(-1)$-无上同调}（或\Term{universally-locally-acyclic}{普适局部 $(-1)$-无上同调}）的\Term{morphism}{态射}，是否都是可变\Term{prescheme}{预概形}上的\Term{etale}{平展}\Term{sheaf}{层}的\Term{fibred-category}{纤维范畴}的\Term{effective-descent}{有效下降}\Term{morphism}{态射}。此陈述似乎颇可信，并将加强 4.3 及第 VIII 讲 9.4
\SourcePage{140}
c)、d)。可将它与 Murre，\textit{Sém. Bourbaki}，第 $293$ 讲，第 $17$ 页的\Term{effective}{有效}性陈述作比较。

c) 在 4.2 的条件下，$g$ 很可能甚至对 $L$ \Term{locally-acyclic}{局部无上同调}且对 $L$ \Term{locally-one-aspherical}{局部 $1$-非球面}，其中 $L$ 为不同于 $Y$ 的各\Term{residue-characteristic}{剩余特征}的素数所构成的\Term{set}{集合}。在此问题中，显然可假设 $Y$ 为\Term{algebraically-closed-field}{代数闭域} $k$ 的谱；可以证明（SGA 1964/65），当 $k$ 上\Term{finite-type}{有限型}\Term{scheme}{概形}可作\Term{resolution-of-singularities}{奇点消解}时，答案为肯定。因此，由 Hironaka 的结果，当 $Y$ 的\Term{characteristic}{特征}为零时，答案为肯定。\footnote{参见 SGA 1 第 XIII 讲。}
\SourcePage{141}
